% Options for packages loaded elsewhere
% Options for packages loaded elsewhere
\PassOptionsToPackage{unicode}{hyperref}
\PassOptionsToPackage{hyphens}{url}
\PassOptionsToPackage{dvipsnames,svgnames,x11names}{xcolor}
%
\documentclass[
  letterpaper,
  DIV=11,
  numbers=noendperiod]{scrartcl}
\usepackage{xcolor}
\usepackage{amsmath,amssymb}
\setcounter{secnumdepth}{-\maxdimen} % remove section numbering
\usepackage{iftex}
\ifPDFTeX
  \usepackage[T1]{fontenc}
  \usepackage[utf8]{inputenc}
  \usepackage{textcomp} % provide euro and other symbols
\else % if luatex or xetex
  \usepackage{unicode-math} % this also loads fontspec
  \defaultfontfeatures{Scale=MatchLowercase}
  \defaultfontfeatures[\rmfamily]{Ligatures=TeX,Scale=1}
\fi
\usepackage{lmodern}
\ifPDFTeX\else
  % xetex/luatex font selection
\fi
% Use upquote if available, for straight quotes in verbatim environments
\IfFileExists{upquote.sty}{\usepackage{upquote}}{}
\IfFileExists{microtype.sty}{% use microtype if available
  \usepackage[]{microtype}
  \UseMicrotypeSet[protrusion]{basicmath} % disable protrusion for tt fonts
}{}
\makeatletter
\@ifundefined{KOMAClassName}{% if non-KOMA class
  \IfFileExists{parskip.sty}{%
    \usepackage{parskip}
  }{% else
    \setlength{\parindent}{0pt}
    \setlength{\parskip}{6pt plus 2pt minus 1pt}}
}{% if KOMA class
  \KOMAoptions{parskip=half}}
\makeatother
% Make \paragraph and \subparagraph free-standing
\makeatletter
\ifx\paragraph\undefined\else
  \let\oldparagraph\paragraph
  \renewcommand{\paragraph}{
    \@ifstar
      \xxxParagraphStar
      \xxxParagraphNoStar
  }
  \newcommand{\xxxParagraphStar}[1]{\oldparagraph*{#1}\mbox{}}
  \newcommand{\xxxParagraphNoStar}[1]{\oldparagraph{#1}\mbox{}}
\fi
\ifx\subparagraph\undefined\else
  \let\oldsubparagraph\subparagraph
  \renewcommand{\subparagraph}{
    \@ifstar
      \xxxSubParagraphStar
      \xxxSubParagraphNoStar
  }
  \newcommand{\xxxSubParagraphStar}[1]{\oldsubparagraph*{#1}\mbox{}}
  \newcommand{\xxxSubParagraphNoStar}[1]{\oldsubparagraph{#1}\mbox{}}
\fi
\makeatother

\usepackage{color}
\usepackage{fancyvrb}
\newcommand{\VerbBar}{|}
\newcommand{\VERB}{\Verb[commandchars=\\\{\}]}
\DefineVerbatimEnvironment{Highlighting}{Verbatim}{commandchars=\\\{\}}
% Add ',fontsize=\small' for more characters per line
\usepackage{framed}
\definecolor{shadecolor}{RGB}{241,243,245}
\newenvironment{Shaded}{\begin{snugshade}}{\end{snugshade}}
\newcommand{\AlertTok}[1]{\textcolor[rgb]{0.68,0.00,0.00}{#1}}
\newcommand{\AnnotationTok}[1]{\textcolor[rgb]{0.37,0.37,0.37}{#1}}
\newcommand{\AttributeTok}[1]{\textcolor[rgb]{0.40,0.45,0.13}{#1}}
\newcommand{\BaseNTok}[1]{\textcolor[rgb]{0.68,0.00,0.00}{#1}}
\newcommand{\BuiltInTok}[1]{\textcolor[rgb]{0.00,0.23,0.31}{#1}}
\newcommand{\CharTok}[1]{\textcolor[rgb]{0.13,0.47,0.30}{#1}}
\newcommand{\CommentTok}[1]{\textcolor[rgb]{0.37,0.37,0.37}{#1}}
\newcommand{\CommentVarTok}[1]{\textcolor[rgb]{0.37,0.37,0.37}{\textit{#1}}}
\newcommand{\ConstantTok}[1]{\textcolor[rgb]{0.56,0.35,0.01}{#1}}
\newcommand{\ControlFlowTok}[1]{\textcolor[rgb]{0.00,0.23,0.31}{\textbf{#1}}}
\newcommand{\DataTypeTok}[1]{\textcolor[rgb]{0.68,0.00,0.00}{#1}}
\newcommand{\DecValTok}[1]{\textcolor[rgb]{0.68,0.00,0.00}{#1}}
\newcommand{\DocumentationTok}[1]{\textcolor[rgb]{0.37,0.37,0.37}{\textit{#1}}}
\newcommand{\ErrorTok}[1]{\textcolor[rgb]{0.68,0.00,0.00}{#1}}
\newcommand{\ExtensionTok}[1]{\textcolor[rgb]{0.00,0.23,0.31}{#1}}
\newcommand{\FloatTok}[1]{\textcolor[rgb]{0.68,0.00,0.00}{#1}}
\newcommand{\FunctionTok}[1]{\textcolor[rgb]{0.28,0.35,0.67}{#1}}
\newcommand{\ImportTok}[1]{\textcolor[rgb]{0.00,0.46,0.62}{#1}}
\newcommand{\InformationTok}[1]{\textcolor[rgb]{0.37,0.37,0.37}{#1}}
\newcommand{\KeywordTok}[1]{\textcolor[rgb]{0.00,0.23,0.31}{\textbf{#1}}}
\newcommand{\NormalTok}[1]{\textcolor[rgb]{0.00,0.23,0.31}{#1}}
\newcommand{\OperatorTok}[1]{\textcolor[rgb]{0.37,0.37,0.37}{#1}}
\newcommand{\OtherTok}[1]{\textcolor[rgb]{0.00,0.23,0.31}{#1}}
\newcommand{\PreprocessorTok}[1]{\textcolor[rgb]{0.68,0.00,0.00}{#1}}
\newcommand{\RegionMarkerTok}[1]{\textcolor[rgb]{0.00,0.23,0.31}{#1}}
\newcommand{\SpecialCharTok}[1]{\textcolor[rgb]{0.37,0.37,0.37}{#1}}
\newcommand{\SpecialStringTok}[1]{\textcolor[rgb]{0.13,0.47,0.30}{#1}}
\newcommand{\StringTok}[1]{\textcolor[rgb]{0.13,0.47,0.30}{#1}}
\newcommand{\VariableTok}[1]{\textcolor[rgb]{0.07,0.07,0.07}{#1}}
\newcommand{\VerbatimStringTok}[1]{\textcolor[rgb]{0.13,0.47,0.30}{#1}}
\newcommand{\WarningTok}[1]{\textcolor[rgb]{0.37,0.37,0.37}{\textit{#1}}}

\usepackage{longtable,booktabs,array}
\usepackage{calc} % for calculating minipage widths
% Correct order of tables after \paragraph or \subparagraph
\usepackage{etoolbox}
\makeatletter
\patchcmd\longtable{\par}{\if@noskipsec\mbox{}\fi\par}{}{}
\makeatother
% Allow footnotes in longtable head/foot
\IfFileExists{footnotehyper.sty}{\usepackage{footnotehyper}}{\usepackage{footnote}}
\makesavenoteenv{longtable}
\usepackage{graphicx}
\makeatletter
\newsavebox\pandoc@box
\newcommand*\pandocbounded[1]{% scales image to fit in text height/width
  \sbox\pandoc@box{#1}%
  \Gscale@div\@tempa{\textheight}{\dimexpr\ht\pandoc@box+\dp\pandoc@box\relax}%
  \Gscale@div\@tempb{\linewidth}{\wd\pandoc@box}%
  \ifdim\@tempb\p@<\@tempa\p@\let\@tempa\@tempb\fi% select the smaller of both
  \ifdim\@tempa\p@<\p@\scalebox{\@tempa}{\usebox\pandoc@box}%
  \else\usebox{\pandoc@box}%
  \fi%
}
% Set default figure placement to htbp
\def\fps@figure{htbp}
\makeatother





\setlength{\emergencystretch}{3em} % prevent overfull lines

\providecommand{\tightlist}{%
  \setlength{\itemsep}{0pt}\setlength{\parskip}{0pt}}



 


\usepackage{mathpazo}
\usepackage{booktabs}
\usepackage{rotating}
\usepackage{makecell}
%% \usepackage[a4paper,margin=2cm]{geometry}

\usepackage{pdfpages}
\usepackage{xcolor}  % For colored warning boxes
\usepackage{pifont}  % For dingbat symbols in error boxes
\usepackage{fontspec}  % For loading the ORCID icon font directly by file
\usepackage[font=Large,labelfont=bf,textfont=bf]{caption}
\DeclareCaptionLabelFormat{addC}{#1 C#2}

%% Load hyperref after other packages to avoid conflicts
\usepackage{hyperref}
%% Set hyperlinks in colour to the "CODECHECK colour".
\hypersetup{colorlinks=true, urlcolor = [rgb]{0,.5,0.2}}

%% ORCID icon as clickable link. Loaded directly from the academicons TTF
%% (shipped with the academicons LaTeX package) by filename via \newfontface,
%% rather than through the academicons package's \aiOrcid command. \aiOrcid
%% selects the font by OS-level family name via fontspec's \newfontfamily,
%% which silently falls back to a legacy 8-bit TFM with no ORCID glyph on
%% systems where the font isn't registered with the OS font system
%% (fontconfig) - producing no visible icon with no compile error. Loading by
%% filename resolves the font via TeX's own search path instead, so it works
%% regardless of OS font registration.
\newfontface\orcidfont{academicons.ttf}
\definecolor{orcidlogocol}{HTML}{A6CE39}
\newcommand\orcidicon[1]{\href{https://orcid.org/#1}{\textcolor{orcidlogocol}{\orcidfont\symbol{"E9D9}}}}
\KOMAoption{captions}{tableheading}
\makeatletter
\@ifpackageloaded{caption}{}{\usepackage{caption}}
\AtBeginDocument{%
\ifdefined\contentsname
  \renewcommand*\contentsname{Table of contents}
\else
  \newcommand\contentsname{Table of contents}
\fi
\ifdefined\listfigurename
  \renewcommand*\listfigurename{List of Figures}
\else
  \newcommand\listfigurename{List of Figures}
\fi
\ifdefined\listtablename
  \renewcommand*\listtablename{List of Tables}
\else
  \newcommand\listtablename{List of Tables}
\fi
\ifdefined\figurename
  \renewcommand*\figurename{Figure}
\else
  \newcommand\figurename{Figure}
\fi
\ifdefined\tablename
  \renewcommand*\tablename{Table}
\else
  \newcommand\tablename{Table}
\fi
}
\@ifpackageloaded{float}{}{\usepackage{float}}
\floatstyle{ruled}
\@ifundefined{c@chapter}{\newfloat{codelisting}{h}{lop}}{\newfloat{codelisting}{h}{lop}[chapter]}
\floatname{codelisting}{Listing}
\newcommand*\listoflistings{\listof{codelisting}{List of Listings}}
\makeatother
\makeatletter
\makeatother
\makeatletter
\@ifpackageloaded{caption}{}{\usepackage{caption}}
\@ifpackageloaded{subcaption}{}{\usepackage{subcaption}}
\makeatother
\usepackage{bookmark}
\IfFileExists{xurl.sty}{\usepackage{xurl}}{} % add URL line breaks if available
\urlstyle{same}
\hypersetup{
  colorlinks=true,
  linkcolor={blue},
  filecolor={Maroon},
  citecolor={Blue},
  urlcolor={Blue},
  pdfcreator={LaTeX via pandoc}}


\author{}
\date{}
\begin{document}


\title{CODECHECK certificate 2026-024}
\maketitle

\centerline{\includegraphics[width=4cm]{/home/schrist/.asdf/installs/r/4.3.2/lib/R/library/codecheck/extdata/codecheck_logo.pdf}}\vspace*{1cm}

\begin{center}
\fcolorbox{red}{yellow}{\parbox{0.9\textwidth}{\centering
\textbf{\Large \textcolor{red}{⚠} WARNING \textcolor{red}{⚠}}\\
\vspace{0.2cm}
\textbf{Report DOI is a placeholder: \texttt{NOT SET}}\\
\vspace{0.1cm}
This certificate is not yet finalized.\\Please set valid identifiers before publishing.
}}
\end{center}

\begin{verbatim}
Warning: Report DOI 'NOT SET' is a placeholder. Please set
valid values before finalizing.
\end{verbatim}

\begin{table}[ht]
\centering
\begin{tabular}{lp{10cm}}
  \hline
Item & Value \\ 
  \hline
Title of checked publication & Beyond colour gamuts: Novel metrics for the reproduction of photoreceptor signals \\ 
  Author(s) & A. C. Hexley \orcidicon{0000-0002-0173-540X} , T. Morimoto \orcidicon{0000-0001-9897-8338} , H. E. Smithson \orcidicon{0000-0001-8904-8315} , M. Spitschan \orcidicon{0000-0002-8572-9268}  \\ 
  Reference & \url{https://doi.org/10.1101/2021.02.27.433203} \\ 
  Codechecker(s) & Simon Christ \orcidicon{0000-0002-5866-1472}  \\ 
  Date of check &  \\ 
  Summary & The code check was straightforward. Only the toolboxes "Statistics and Machine Learning" and "Image Processing" needed to be added in contrast to what was stated in the README. \\ 
  Repository & \url{https://github.com/takuma929/metrics_photoreceptor_signal_reproduction} \\ 
   \hline
\end{tabular}
\caption{CODECHECK summary} 
\end{table}

\begin{table}[ht]
\centering
\begin{tabular}{p{6cm}p{6cm}p{2cm}}
  \hline
Output & Comment & Size (b) \\ 
  \hline
\href{https://github.com/takuma929/metrics_photoreceptor_signal_reproduction/blob/master/codecheck/outputs/fig1a.pdf}{\path{figs/fig1a.pdf}} & Fig 1 - spectral sensitivities & example signals (panels fig1a-fig1h) &  \\ 
  \href{https://github.com/takuma929/metrics_photoreceptor_signal_reproduction/blob/master/codecheck/outputs/fig2a.pdf}{\path{figs/fig2a.pdf}} & Fig 2 - equal-luminance MacLeod-Boynton locus (panels fig2a-fig2c + inserts) &  \\ 
  \href{https://github.com/takuma929/metrics_photoreceptor_signal_reproduction/blob/master/codecheck/outputs/fig3a.pdf}{\path{figs/fig3a.pdf}} & Fig 3 - display primaries (panels fig3a-fig3h) &  \\ 
  \href{https://github.com/takuma929/metrics_photoreceptor_signal_reproduction/blob/master/codecheck/outputs/fig4a.png}{\path{figs/fig4a.png}} & Fig 4a - chromaticity reproduction diagram &  \\ 
  \href{https://github.com/takuma929/metrics_photoreceptor_signal_reproduction/blob/master/codecheck/outputs/fig4b.pdf}{\path{figs/fig4b.pdf}} & Fig 4b - chromaticity reproduction bar graph &  \\ 
  \href{https://github.com/takuma929/metrics_photoreceptor_signal_reproduction/blob/master/codecheck/outputs/fig4c.pdf}{\path{figs/fig4c.pdf}} & Fig 4c - PSRM (photoreceptor signal reproduction) bar graph &  \\ 
  \href{https://github.com/takuma929/metrics_photoreceptor_signal_reproduction/blob/master/codecheck/outputs/fig5ai.pdf}{\path{figs/fig5ai.pdf}} & Fig 5 - PSDM rod & melanopsin distortion (panels fig5ai-fig5cii) &  \\ 
  \href{https://github.com/takuma929/metrics_photoreceptor_signal_reproduction/blob/master/codecheck/outputs/fig6a_box_CRT.pdf}{\path{figs/fig6a_box_CRT.pdf}} & Fig 6 - PSDM per-display box/xy plots (CRT, LCD, DP variants) &  \\ 
  \href{https://github.com/takuma929/metrics_photoreceptor_signal_reproduction/blob/master/codecheck/outputs/fig7_LvsS_all.pdf}{\path{figs/fig7_LvsS_all.pdf}} & Fig 7 - equal-luminance excitation diagrams, full database + 3-primary displays (fig7_{LvsS,LvsI,IvsS}_{all,crt,lcd,dp}) &  \\ 
  \href{https://github.com/takuma929/metrics_photoreceptor_signal_reproduction/blob/master/codecheck/outputs/fig8_LvsS_MPHDR.pdf}{\path{figs/fig8_LvsS_MPHDR.pdf}} & Fig 8 - equal-luminance excitation diagrams, 5-primary displays (fig8_{LvsS,LvsI,IvsS}_{MPHDR,Man,nb5p,bb5p,NZ}) &  \\ 
  \href{https://github.com/takuma929/metrics_photoreceptor_signal_reproduction/blob/master/codecheck/outputs/fig9a.pdf}{\path{figs/fig9a.pdf}} & Fig 9 - bit depth & tolerance analysis (fig9a, fig9b) &  \\ 
   \hline
\end{tabular}
\caption{Summary of output files generated} 
\end{table}

\section{Summary}\label{summary}

This code was (not) straightforward to codecheck. {[}\ldots{} ADD MORE
INTERESTING FINDINGS HERE..{]}

\clearpage

\section{CODECHECKER notes}\label{codechecker-notes}

The GitHub repo \ldots{} This check is based on the commit
\texttt{d1fcfb87a67a35a62011b873623d81abf5286c0f}. Code was written in
\ldots{} I went through the following steps \ldots{} One hard problem
was \ldots{} I added the following files \ldots{} using tools/methods
\ldots{}

This took \ldots{} minutes to complete on \{a large workstation, my
laptop\}.

\subsection{Recommendations}\label{recommendations}

I suggest to the authors to consider the following suggestions for their
next publication or workflow:

\begin{itemize}
\tightlist
\item
  \ldots{}
\end{itemize}

(document here if any of the suggestions were taken up by the authors in
the meantime - do not remove any, keep track of contributions via
feedback)

\clearpage

\section{Manifest files}\label{manifest-files}

\subsection{fig1a.pdf}\label{fig1a.pdf}

\textbf{Comment:} Fig 1 - spectral sensitivities \& example signals
(panels fig1a-fig1h)

\pandocbounded{\includegraphics[keepaspectratio]{./home/schrist/projects/SchneiderMorimotoEtAl_ColorResAppl_2026/codecheck/outputs/fig1a.pdf}}
\clearpage 

\subsection{fig2a.pdf}\label{fig2a.pdf}

\textbf{Comment:} Fig 2 - equal-luminance MacLeod-Boynton locus (panels
fig2a-fig2c + inserts)

\pandocbounded{\includegraphics[keepaspectratio]{./home/schrist/projects/SchneiderMorimotoEtAl_ColorResAppl_2026/codecheck/outputs/fig2a.pdf}}
\clearpage 

\subsection{fig3a.pdf}\label{fig3a.pdf}

\textbf{Comment:} Fig 3 - display primaries (panels fig3a-fig3h)

\pandocbounded{\includegraphics[keepaspectratio]{./home/schrist/projects/SchneiderMorimotoEtAl_ColorResAppl_2026/codecheck/outputs/fig3a.pdf}}
\clearpage 

\subsection{fig4a.png}\label{fig4a.png}

\textbf{Comment:} Fig 4a - chromaticity reproduction diagram

\pandocbounded{\includegraphics[keepaspectratio]{./home/schrist/projects/SchneiderMorimotoEtAl_ColorResAppl_2026/codecheck/outputs/fig4a.png}}
\clearpage 

\subsection{fig4b.pdf}\label{fig4b.pdf}

\textbf{Comment:} Fig 4b - chromaticity reproduction bar graph

\pandocbounded{\includegraphics[keepaspectratio]{./home/schrist/projects/SchneiderMorimotoEtAl_ColorResAppl_2026/codecheck/outputs/fig4b.pdf}}
\clearpage 

\subsection{fig4c.pdf}\label{fig4c.pdf}

\textbf{Comment:} Fig 4c - PSRM (photoreceptor signal reproduction) bar
graph

\pandocbounded{\includegraphics[keepaspectratio]{./home/schrist/projects/SchneiderMorimotoEtAl_ColorResAppl_2026/codecheck/outputs/fig4c.pdf}}
\clearpage 

\subsection{fig5ai.pdf}\label{fig5ai.pdf}

\textbf{Comment:} Fig 5 - PSDM rod \& melanopsin distortion (panels
fig5ai-fig5cii)

\pandocbounded{\includegraphics[keepaspectratio]{./home/schrist/projects/SchneiderMorimotoEtAl_ColorResAppl_2026/codecheck/outputs/fig5ai.pdf}}
\clearpage 

\subsection{fig6a\_box\_CRT.pdf}\label{fig6a_box_crt.pdf}

\textbf{Comment:} Fig 6 - PSDM per-display box/xy plots (CRT, LCD, DP
variants)

\pandocbounded{\includegraphics[keepaspectratio]{./home/schrist/projects/SchneiderMorimotoEtAl_ColorResAppl_2026/codecheck/outputs/fig6a_box_CRT.pdf}}
\clearpage 

\subsection{fig7\_LvsS\_all.pdf}\label{fig7_lvss_all.pdf}

\textbf{Comment:} Fig 7 - equal-luminance excitation diagrams, full
database + 3-primary displays
(fig7\_\{LvsS,LvsI,IvsS\}\_\{all,crt,lcd,dp\})

\pandocbounded{\includegraphics[keepaspectratio]{./home/schrist/projects/SchneiderMorimotoEtAl_ColorResAppl_2026/codecheck/outputs/fig7_LvsS_all.pdf}}
\clearpage 

\subsection{fig8\_LvsS\_MPHDR.pdf}\label{fig8_lvss_mphdr.pdf}

\textbf{Comment:} Fig 8 - equal-luminance excitation diagrams, 5-primary
displays (fig8\_\{LvsS,LvsI,IvsS\}\_\{MPHDR,Man,nb5p,bb5p,NZ\})

\pandocbounded{\includegraphics[keepaspectratio]{./home/schrist/projects/SchneiderMorimotoEtAl_ColorResAppl_2026/codecheck/outputs/fig8_LvsS_MPHDR.pdf}}
\clearpage 

\subsection{fig9a.pdf}\label{fig9a.pdf}

\textbf{Comment:} Fig 9 - bit depth \& tolerance analysis (fig9a, fig9b)

\pandocbounded{\includegraphics[keepaspectratio]{./home/schrist/projects/SchneiderMorimotoEtAl_ColorResAppl_2026/codecheck/outputs/fig9a.pdf}}
\clearpage 

\clearpage

\subsection{Acknowledgements}\label{acknowledgements}

I would like to thank \ldots{} for \ldots{} {[}promptly answering any
queries/providing updates; assitance with X; \ldots{]}.

\section{Citing this document}\label{citing-this-document}

\section{About CODECHECK}\label{about-codecheck}

This certificate confirms that the codechecker could independently
reproduce the results of a computational analysis given the data and
code from a third party. A CODECHECK does not check whether the original
computation analysis is correct. However, as all materials required for
the reproduction are freely available by following the links in this
document, the reader can then study for themselves the code and data.

\subsection{Multi-language code
chunks}\label{multi-language-code-chunks}

Quarto can execute R, Python, and Julia code in the same document - see
the \href{https://quarto.org/docs/computations/languages.html}{Quarto
language guides} for details. The R chunk below is a real, unevaluated
(\texttt{eval:\ false}) knitr chunk. The Python and Julia snippets are
shown as plain (non-executed) code blocks rather than
\texttt{\{python\}}/\texttt{\{julia\}} engine chunks, since those
engines need the \texttt{reticulate} or \texttt{JuliaCall} packages (and
a Python/Julia installation) even to be parsed - not needed just to
demonstrate the syntax. Use \texttt{\{python\}}/\texttt{\{julia\}}
chunks instead if your report actually executes Python or Julia code.

\begin{Shaded}
\begin{Highlighting}[]
\DecValTok{1} \SpecialCharTok{+} \DecValTok{1}
\end{Highlighting}
\end{Shaded}

\begin{Shaded}
\begin{Highlighting}[]
\DecValTok{1} \OperatorTok{+} \DecValTok{1}
\end{Highlighting}
\end{Shaded}

\begin{Shaded}
\begin{Highlighting}[]
\FloatTok{1} \OperatorTok{+} \FloatTok{1}
\end{Highlighting}
\end{Shaded}

\section{About this document}\label{about-this-document}

This document was created using \href{https://quarto.org/}{Quarto} using
the \href{https://github.com/codecheckers/codecheck}{\texttt{codecheck}}
R package. See the
\href{https://quarto.org/docs/output-formats/pdf-basics.html}{Quarto PDF
basics} and
\href{https://quarto.org/docs/output-formats/pdf-engine.html}{PDF
engine} documentation for background on the PDF output produced here
(via tinytex). \texttt{make\ codecheck.pdf} will regenerate the report
file.

\begin{Shaded}
\begin{Highlighting}[]
\FunctionTok{sessionInfo}\NormalTok{()}
\end{Highlighting}
\end{Shaded}

\begin{verbatim}
R version 4.3.2 (2023-10-31)
Platform: x86_64-pc-linux-gnu (64-bit)
Running under: Ubuntu 22.04.5 LTS

Matrix products: default
BLAS:   /home/schrist/.asdf/installs/r/4.3.2/lib/R/lib/libRblas.so 
LAPACK: /usr/lib/x86_64-linux-gnu/openblas-pthread/liblapack.so.3;  LAPACK version 3.10.0

locale:
 [1] LC_CTYPE=en_US.UTF-8       LC_NUMERIC=C              
 [3] LC_TIME=de_DE.UTF-8        LC_COLLATE=en_US.UTF-8    
 [5] LC_MONETARY=de_DE.UTF-8    LC_MESSAGES=en_US.UTF-8   
 [7] LC_PAPER=de_DE.UTF-8       LC_NAME=C                 
 [9] LC_ADDRESS=C               LC_TELEPHONE=C            
[11] LC_MEASUREMENT=de_DE.UTF-8 LC_IDENTIFICATION=C       

time zone: Europe/Berlin
tzcode source: system (glibc)

attached base packages:
[1] stats     graphics  grDevices utils     datasets 
[6] methods   base     

other attached packages:
[1] tibble_3.3.1          xtable_1.8-8         
[3] yaml_2.3.12           rprojroot_2.1.1      
[5] knitr_1.52            codecheck_0.30.0.9000
[7] parsedate_1.3.2       R.cache_0.17.0       
[9] gh_1.6.1             

loaded via a namespace (and not attached):
 [1] utf8_1.2.6        generics_0.1.4    xml2_1.6.0       
 [4] zen4R_0.10.6      stringi_1.8.9     httpcode_0.3.0   
 [7] digest_0.6.39     magrittr_2.0.5    evaluate_1.0.5   
[10] fastmap_1.2.0     plyr_1.8.9        cellranger_1.1.0 
[13] R.oo_1.27.1       jsonlite_2.0.0    R.utils_2.13.0   
[16] whisker_0.4.1     crul_1.6.0        httr_1.4.9       
[19] purrr_1.2.2       XML_3.99-0.24     httr2_1.3.0      
[22] cli_3.6.6         pdftools_3.9.1    rlang_1.3.0      
[25] R.methodsS3_1.8.2 gitcreds_0.1.2    withr_3.0.3      
[28] cachem_1.1.0      otel_0.2.0        tools_4.3.2      
[31] memoise_2.0.1     dplyr_1.2.1       assertthat_0.2.1 
[34] curl_8.0.0        vctrs_0.7.3       R6_2.6.1         
[37] lifecycle_1.0.5   magick_2.9.1      git2r_0.36.2     
[40] stringr_1.6.0     rsvg_2.7.0        fs_2.1.0         
[43] pkgconfig_2.0.3   rorcid_0.7.0      osfr_0.2.9       
[46] pillar_1.11.1     glue_1.8.1        Rcpp_1.1.2       
[49] xfun_0.61         tidyselect_1.2.1  keyring_1.4.1    
[52] htmltools_0.5.9   rmarkdown_2.32    qpdf_1.4.1       
[55] compiler_4.3.2    askpass_1.2.1     fauxpas_0.6.0    
[58] readxl_1.5.0.1   
\end{verbatim}




\end{document}
